% ═══════════════════════════════════════════════════════════════
% QR-Code: Lernpfad Verkehrssicherheit
% Zum Projizieren oder Aushängen im Klassenraum
% ═══════════════════════════════════════════════════════════════

\documentclass[a4paper, 11pt]{article}

\usepackage[utf8]{inputenc}
\usepackage[T1]{fontenc}
\usepackage[ngerman]{babel}
\usepackage{helvet}
\renewcommand{\familydefault}{\sfdefault}

\usepackage[margin=2cm]{geometry}
\usepackage{xcolor}
\usepackage{qrcode}
\usepackage[hidelinks]{hyperref}

% Farben
\definecolor{physik}{HTML}{2c5aa0}
\definecolor{physik-hell}{HTML}{e8f0fa}

% URL
\newcommand{\lpurl}{https://devbox42.github.io/sciencesim/physik/kl06/02-bewegung/std-04-verkehrssicherheit/LP-04-verkehrssicherheit.html}

\setlength{\parindent}{0pt}

\begin{document}

\pagestyle{empty}

\begin{center}

\vspace*{1cm}

{\Huge\bfseries\textcolor{physik}{Lernpfad}}

\vspace{0.3cm}

{\LARGE Verkehrssicherheit: Anhalteweg}

\vspace{0.2cm}

{\large Physik Klasse 6 | Bewegungen}

\vspace{1cm}

\textcolor{physik}{\rule{12cm}{2pt}}

\vspace{1.5cm}

% QR-Code (groß, klickbar)
\href{\lpurl}{\qrcode[height=8cm]{\lpurl}}

\vspace{1.5cm}

\textcolor{physik}{\rule{12cm}{2pt}}

\vspace{0.8cm}

{\large Scanne den QR-Code mit deinem Tablet oder Smartphone!}

\vspace{0.5cm}

{\small\textcolor{gray}{\url{\lpurl}}}

\vspace{1cm}

\fcolorbox{physik}{physik-hell}{%
    \parbox{10cm}{%
        \centering
        \vspace{0.3cm}
        \textbf{Inhalt:}\\[0.2em]
        Schrecksekunde $\bullet$ Reaktionsweg $\bullet$ Bremsweg\\[0.1em]
        Anhalteweg $\bullet$ Simulation $\bullet$ Mini-Quiz
        \vspace{0.3cm}
    }%
}

\end{center}

\end{document}
